\clearpage
\section*{译注汇总}
\addcontentsline{toc}{section}{译注汇总}

本附录把正文中的全部译注按出现顺序重列一遍，并注明所在节、所对照的 nLab 词条与来源等级，以便不翻正文即可通读批注。正文中的译注为页脚小字，此处为便于检索作了压缩改写（不是逐字复制），语义以正文脚注为准。

\medskip
\noindent\textbf{引言　GLC 的两个阶段（A 级）}\quad 作者把几何朗兰兹的发展分为两期：第一阶段（Phase I）自问题的提出（[Dri83, Lau87]）延伸到 $\mathrm{GL}_{n}$ 的 Hecke 本征层的构造（[FGV01, Gai04]），早期综述见 [Fre04, Gai17a]；第二阶段（Phase II）即范畴化的 GLC（categorical GLC），以 [BD] 为起点，经 Arinkin--Gaitsgory 以奇异支撑与 $\IndCoh_{\Nilp}$ 重述（[AG15]），直至 2024 年非分歧猜想的证明（[GR24a, GR24b, ABC${}^{+}$24a, ABC${}^{+}$24b, CCF${}^{+}$24]）。分期系作者自述，nLab 无对应词条，故不附链接。

\smallskip
\noindent\textbf{引言　Current Events Bulletin（B 级）}\quad 美国数学会在年度联合数学会议（Joint Mathematics Meetings）上组织的综述报告系列，本文即该系列报告之一，致谢对象为当届组织者 Daniel Erman 与 Bianca Viray。见 \url{https://www.ams.org/meetings/lectures/current-events-bulletin}（2026 年 11 月可访问）。

\smallskip
\noindent\textbf{§1.1　谱分解与 Tannaka 对偶（B 级）}\quad 把谱定理读作“表示范畴 $\simeq$ 谱侧范畴”，实质是 Tannaka 对偶在交换群上的特例：由表示范畴连同遗忘函子可重构出群，谱侧即“颜色／参数”的栖身之处。对照词条：\emph{Tannaka duality}（\href{https://ncatlab.org/nlab/show/Tannaka+duality}{\texttt{https:/\allowbreak{}/\allowbreak{}ncatlab.\allowbreak{}org/\allowbreak{}nlab/\allowbreak{}show/\allowbreak{}Tannaka+\allowbreak{}duality}}）。这一读法预示了 GLC 的基本形状：一侧是范畴，另一侧是其谱。

\smallskip
\noindent\textbf{§1.4　朗兰兹对偶群（B 级）}\quad $\Gd$ 的组合定义只用到根数据（root datum）的对偶：保留外尔群，把环面与根空间的结构对偶化。对照词条：\emph{Langlands dual group}（\href{https://ncatlab.org/nlab/show/Langlands+dual+group}{\texttt{https:/\allowbreak{}/\allowbreak{}ncatlab.\allowbreak{}org/\allowbreak{}nlab/\allowbreak{}show/\allowbreak{}Langlands+\allowbreak{}dual+\allowbreak{}group}}），其中有 $\mathrm{SL}_{n}\leftrightarrow\mathrm{PGL}_{n}$、$\mathrm{SO}_{2n+1}\leftrightarrow\mathrm{Sp}_{2n}$ 等例。本译本统一作“朗兰兹对偶群”，其后简称“对偶群”。

\smallskip
\noindent\textbf{§1.4　S-对偶与 Kapustin--Witten（B 级）}\quad “电磁对偶把 $G$ 换成 $\Gd$”在四维 $\mathcal{N}=4$ 超对称 Yang--Mills 中由 Montonen--Olive 对偶实现；Kapustin--Witten 把它进一步化为几何朗兰兹对偶。对照词条：\emph{S-duality}（\href{https://ncatlab.org/nlab/show/S-duality}{\texttt{https:/\allowbreak{}/\allowbreak{}ncatlab.\allowbreak{}org/\allowbreak{}nlab/\allowbreak{}show/\allowbreak{}S-duality}}）、\emph{Kapustin-Witten TQFT}（\href{https://ncatlab.org/nlab/show/Kapustin-Witten+TQFT}{\texttt{https:/\allowbreak{}/\allowbreak{}ncatlab.\allowbreak{}org/\allowbreak{}nlab/\allowbreak{}show/\allowbreak{}Kapustin-Witten+\allowbreak{}TQFT}}）。

\smallskip
\noindent\textbf{§2.1　六函子形式（B 级）}\quad 指把层的推前、拉回、张量积、内 Hom、上同调、紧支撑上同调六个操作在导出范畴层面公理化并使之相容的形式体系。原文“brave new harmonic analysis”的要旨是：泛函分析里的积分、乘法、函数空间，在层论中分别对应推前／直积分、张量积、层的导出范畴。对照词条：\emph{six operations}（\href{https://ncatlab.org/nlab/show/six+operations}{\texttt{https:/\allowbreak{}/\allowbreak{}ncatlab.\allowbreak{}org/\allowbreak{}nlab/\allowbreak{}show/\allowbreak{}six+\allowbreak{}operations}}）。

\smallskip
\noindent\textbf{§2.2.2　Weil 单值化（B 级）}\quad 阿代尔描述与 Weil 单值化定理是同一陈述：在有限多个点外把 $G$-丛平凡化，用 $G(F)$ 粘起来，剩余自由度是 $\prod'_{x}G(\C(\!(z_{x})\!))$ 模 $\prod_{x}G(\C[[z_{x}]])$。对照词条：\emph{Weil uniformization theorem}（\href{https://ncatlab.org/nlab/show/Weil+uniformization+theorem}{\texttt{https:/\allowbreak{}/\allowbreak{}ncatlab.\allowbreak{}org/\allowbreak{}nlab/\allowbreak{}show/\allowbreak{}Weil+\allowbreak{}uniformization+\allowbreak{}theorem}}）。

\smallskip
\noindent\textbf{§2.4　几何类域论（B 级）}\quad $G=\mathrm{GL}_{1}$ 的 GLC，其经典形式即 Abel--Jacobi 给出的 $\mathrm{Pic}^{0}(C)\cong\Loc_{\mathrm{GL}_{1}}(C)$；nLab 把它表述为“曲线基本群的阿贝尔化与其 Picard 簇基本群之等同”，并视之为高维几何 Langlands 的出发点。对照词条：\emph{geometric class field theory}（\href{https://ncatlab.org/nlab/show/geometric+class+field+theory}{\texttt{https:/\allowbreak{}/\allowbreak{}ncatlab.\allowbreak{}org/\allowbreak{}nlab/\allowbreak{}show/\allowbreak{}geometric+\allowbreak{}class+\allowbreak{}field+\allowbreak{}theory}}）。

\smallskip
\noindent\textbf{§2.5　归纳凝聚层（B 级）}\quad $\IndCoh$ 与 $\QCoh$ 之别在于它同时记录推前与反常推前，故能容纳奇异簇上的“错维”层；下标 $\Nilp$ 指幂零奇异支撑（见 §4.2）。对照词条：\emph{ind-coherent sheaves}（\href{https://ncatlab.org/nlab/show/ind-coherent+sheaves}{\texttt{https:/\allowbreak{}/\allowbreak{}ncatlab.\allowbreak{}org/\allowbreak{}nlab/\allowbreak{}show/\allowbreak{}ind-coherent+\allowbreak{}sheaves}}）。

\smallskip
\noindent\textbf{§3.2　因子化同调（B 级）}\quad 把因子化代数在小圆盘上取值、沿嵌入取高阶余极限而得的“整体不变量”；拓扑原型是 Ayala--Françis 对拓扑流形的构造。原文强调的正是“局部数据 $\Rightarrow$ 整体可观测量”的装配机制。对照词条：\emph{factorization homology}（\href{https://ncatlab.org/nlab/show/factorization+homology}{\texttt{https:/\allowbreak{}/\allowbreak{}ncatlab.\allowbreak{}org/\allowbreak{}nlab/\allowbreak{}show/\allowbreak{}factorization+\allowbreak{}homology}}）。

\smallskip
\noindent\textbf{§3.3　几何 Satake 等价（B 级）}\quad $\Sph_{G}$ 是 $\LGp$-等变层构成的交换张量范畴，其 Tannaka 谱恰为 $\Rep(\Gd)$，于是 $\Gd$ 由 $\Sph_{G}$ 的自同构群重建出来——Tannaka 重构的几何版本。对照词条：\emph{geometric Satake equivalence}（\href{https://ncatlab.org/nlab/show/geometric+Satake+equivalence}{\texttt{https:/\allowbreak{}/\allowbreak{}ncatlab.\allowbreak{}org/\allowbreak{}nlab/\allowbreak{}show/\allowbreak{}geometric+\allowbreak{}Satake+\allowbreak{}equivalence}}）、\emph{Tannaka duality}。

\smallskip
\noindent\textbf{§3.5　谱作用（B 级）}\quad 谱侧是 $\QCoh(\Loc_{\Gd})$，“作用”来自各点处几何 Satake 等价与 $\Rep(\Gd)$ 的对称幺半结构，经因子化同调粘成整体；单点情形即几何 Satake。对照词条：\emph{geometric Langlands correspondence}（\href{https://ncatlab.org/nlab/show/geometric+Langlands+correspondence}{\texttt{https:/\allowbreak{}/\allowbreak{}ncatlab.\allowbreak{}org/\allowbreak{}nlab/\allowbreak{}show/\allowbreak{}geometric+\allowbreak{}Langlands+\allowbreak{}correspondence}}）。

\smallskip
\noindent\textbf{§4.2　奇异支撑与微局部（B 级）}\quad 层的奇异支撑是相空间 $T^{*}X$ 中的子簇，记录层在何处、沿何方向“不光滑”，是物理中波前集的代数对应物；对 D-模它与特征簇在 Riemann--Hilbert 对应下相合，$\Nilp$ 正是自守层奇异支撑所在的锥。对照词条：nLab 的 \emph{singular support} 页面实际重定向到“分布的支集”（义项不同），故改引 \emph{characteristic variety}（\href{https://ncatlab.org/nlab/show/characteristic+variety}{\texttt{https:/\allowbreak{}/\allowbreak{}ncatlab.\allowbreak{}org/\allowbreak{}nlab/\allowbreak{}show/\allowbreak{}characteristic+\allowbreak{}variety}}，最后修订 2017-08-29）。

\smallskip
\noindent\textbf{§5.2　周期与相对 Langlands（B 级）}\quad 此处的“周期”不是代数几何的周期积分，而是自守形式上的不变线性泛函（把自守表示“求值／积分”出一个数或向量空间）；相对 Langlands 纲领即系统研究这些周期并与谱侧的“测量”相匹配。对照词条：\emph{relative Langlands program}（\href{https://ncatlab.org/nlab/show/relative+Langlands+program}{\texttt{https:/\allowbreak{}/\allowbreak{}ncatlab.\allowbreak{}org/\allowbreak{}nlab/\allowbreak{}show/\allowbreak{}relative+\allowbreak{}Langlands+\allowbreak{}program}}）。

\subsection*{附：核心概念与 nLab 词条对照表}

下表把文中反复出现、且 nLab 上有专门词条的概念集中列出，供读者按图索骥（均为 B 级来源；可访问性核对于 2026 年 11 月，修订日期取自本机 \texttt{07\_refs\_nlab} 缓存，抓取于 2026-11-15）。nLab 无独立词条的概念（如 opers、分歧情形）此处不作强行对应。

{\small
\begin{longtable}{@{}>{\raggedright\arraybackslash}p{0.26\linewidth}>{\raggedright\arraybackslash}p{0.24\linewidth}>{\raggedright\arraybackslash}p{0.13\linewidth}>{\raggedright\arraybackslash}p{0.26\linewidth}@{}}
\toprule
概念（中／英） & nLab 词条（URL 末段） & 最后修订 & 在文中的位置与作用 \\
\midrule
\endfirsthead
\toprule
概念（中／英） & nLab 词条（URL 末段） & 最后修订 & 在文中的位置与作用 \\
\midrule
\endhead
\bottomrule
\endfoot
几何朗兰兹对偶 & geometric+​\allowbreak Langlands+​\allowbreak correspondence & \mbox{2026-04-11} & 全文主题；§2.5 主定理、§3.5 谱作用 \\
几何 Satake 等价 & geometric+​\allowbreak Satake+​\allowbreak equivalence & \mbox{2024-07-22} & §3.3；对偶群的来源 \\
Tannaka 对偶 & Tannaka+​\allowbreak duality & \mbox{2025-07-28} & §1.1、§3.4；谱重构的语言 \\
朗兰兹对偶群 & Langlands+​\allowbreak dual+​\allowbreak group & \mbox{2014-09-07} & §1.4；$\Gd$ 的定义 \\
六函子形式 & six+​\allowbreak operations & \mbox{2026-09-30} & §2.1；层论的语法 \\
$G$-主丛／向量丛模栈 & moduli+​\allowbreak stack & \mbox{2023-04-16} & §2.2、§2.2.1--2.2.2；$\Bun_{G}$ 与 $\mathrm{Pic}$ 的框架 \\
仿射 Grassmannian & affine+​\allowbreak Grassmannian & \mbox{2023-08-15} & §3.1；Hecke 算子的载体 \\
Weil 单值化 & Weil+​\allowbreak uniformization+​\allowbreak theorem & \mbox{2023-04-04} & §2.2.2；阿代尔描述 \\
几何类域论 & geometric+​\allowbreak class+​\allowbreak field+​\allowbreak theory & \mbox{2025-05-11} & §2.4；$\mathrm{GL}_{1}$ 情形 \\
D-模 & D-module & \mbox{2024-07-27} & §2.4.1；de Rham 侧的基本对象 \\
Riemann--Hilbert 对应 & Riemann-Hilbert+​\allowbreak correspondence & \mbox{2024-08-09} & §2.3；Betti 与 de Rham 的桥 \\
Hecke 本征层 & Hecke+​\allowbreak eigensheaf & — & §2.5、§5.1；对应的“解” \\
归纳凝聚层 & ind-coherent+​\allowbreak sheaves & \mbox{2015-10-05} & §2.5、§5.1；奇异谱侧的修理 \\
因子化代数 & factorization+​\allowbreak algebra & \mbox{2026-08-16} & §3.2；交换性的结构来源 \\
因子化同调 & factorization+​\allowbreak homology & \mbox{2026-05-30} & §3.2；整体装配机制 \\
Hitchin 纤维化 & Hitchin+​\allowbreak fibration & \mbox{2023-11-01} & §4.1；阿贝尔化的几何 \\
Kapustin--Witten TQFT & Kapustin-Witten+​\allowbreak TQFT & \mbox{2026-09-28} & §1.4、§6.2；物理来源 \\
S-对偶 & S-duality & \mbox{2026-09-18} & §1.4、§6.2；$G\leftrightarrow\Gd$ 的物理机制 \\
Fargues--Fontaine 曲线 & Fargues-Fontaine+​\allowbreak curve & \mbox{2022-07-04} & §6.1.1；$p$-进局部 Langlands \\
函数域类比 & function+​\allowbreak field+​\allowbreak analogy & \mbox{2024-06-03} & §6.1.2；Weil 的罗塞塔石碑 \\
量子几何朗兰兹 & quantum+​\allowbreak geometric+​\allowbreak Langlands+​\allowbreak correspondence & \mbox{2026-04-06} & §6.0.4；形变版本 \\
Kac--Moody 代数 & Kac-Moody+​\allowbreak algebra & \mbox{2026-11-06} & §5.5、§6.0.4；opers 的来源 \\
Eisenstein 级数 & Eisenstein+​\allowbreak series & \mbox{2014-03-26} & §5.3；归纳步的连续谱 \\
Whittaker 函数 & Whittaker+​\allowbreak function & \mbox{2023-06-14} & §5.4；白光的归一化 \\
自守形式 & automorphic+​\allowbreak form & \mbox{2026-09-09} & 全文；去范畴化后的对象 \\
Langlands 纲领 & Langlands+​\allowbreak program & \mbox{2024-03-28} & §6.1；算术侧的大背景 \\
反常层 & perverse+​\allowbreak sheaf & \mbox{2026-10-16} & §2.3、§5.1；Betti 侧的自然栖息地 \\
手征代数 & chiral+​\allowbreak algebra & \mbox{2026-11-02} & §3.2、§5.5；因子化代数的场论原型 \\
Beilinson--Drinfeld 代数 & Beilinson-Drinfeld+​\allowbreak algebra & \mbox{2022-10-18} & §3.2；因子化代数的原型 \\
奇异支撑 & characteristic+​\allowbreak variety & \mbox{2017-08-29} & §4.2、§5.1；$\Nilp$ 下标（nLab 的 \emph{singular support} 页实为重定向到“分布的支集”，义项不同） \\
\end{longtable}
}

\noindent\textit{说明：}nLab 为维基式词条，同一概念可能同时存在多个名称（如“opers”尚无独立词条，正文因此只作文字说明而未给链接）。凡 nLab 未设词条者，本译本不以其他网络来源（含各类百科与博客）替代，以符合来源等级纪律。
